\documentclass[11pt, a4paper]{article}

\usepackage[T1]{fontenc}
\usepackage[utf8]{inputenc}
\usepackage{tgpagella}
\usepackage[spanish, es-noshorthands]{babel}
\usepackage{amsmath, amssymb, bm}
\usepackage{booktabs}
\usepackage{tabularx}
\usepackage[table, xcdraw]{xcolor}
\usepackage[most]{tcolorbox}
\usepackage[margin=2.5cm]{geometry}
\usepackage{ragged2e}
\usepackage{fancyhdr} % Agregado para resolver "Undefined control sequence"

\pagestyle{fancy}
\fancyhf{}
\setlength{\headheight}{16pt}
\lhead{\textbf{UCB Sede Tarija} | Ing. Mecatrónica}
\chead{}
\rhead{IMT-121 Circuitos Electrónicos I | Hoja de Datos Lab A}
\lfoot{Prof: M.Sc. Ing. Bernardo Quiroga Turdera}
\rfoot{Página \thepage}
\renewcommand{\headrulewidth}{0.4pt}

\definecolor{ucbblue}{RGB}{0, 51, 102}

\begin{document}

\begin{tcolorbox}[colback=ucbblue!5!white, colframe=ucbblue, title=\textbf{Hoja de Registro de Datos: Experiencia Integrada A}]
    \centering \textbf{Registro de Predicciones y Mediciones en Banco}
\end{tcolorbox}

\vspace{0.2cm}
\textbf{Equipo / Integrantes:} \hrulefill

\section*{Tabla 1: Verificación de Componentes (Fase 1)}
\renewcommand{\arraystretch}{1.5}
% Eliminamos el exceso de ancho forzando la tabla a ser del ancho exacto del texto y sin desbordes.
\noindent
\begin{tabularx}{\textwidth}{|>{\RaggedRight}X|>{\RaggedRight}X|>{\RaggedRight}X|}
\hline
\rowcolor{ucbblue!10}\textbf{Componente} & \textbf{Valor Nominal} & \textbf{Valor Medido ($\Omega$)} \\
\hline
$R_1$ & $4.7\,\text{k}\Omega$ & \\ \hline
$R_2$ & $10\,\text{k}\Omega$ & \\ \hline
$R_3$ & $6.8\,\text{k}\Omega$ & \\ \hline
$R_{L1}$ & $1.0\,\text{k}\Omega$ & \\ \hline
$R_{L2}$ & $2.2\,\text{k}\Omega$ & \\ \hline
$R_{L3}$ & $4.7\,\text{k}\Omega$ & \\ \hline
$R_{L4}$ & $10\,\text{k}\Omega$ & \\ \hline
\end{tabularx}

\section*{Tabla 2: Predicciones Analíticas (Prelaboratorio)}
\noindent
\begin{tabularx}{\textwidth}{|>{\RaggedRight}X|>{\RaggedRight}X|>{\RaggedRight}X|}
\hline
\rowcolor{ucbblue!10}\textbf{Parámetro} & \textbf{Cálculo Analítico} & \textbf{Simulación SPICE} \\
\hline
$V_{a}^{(V1)}$ (Superposición) & & \\ \hline
$V_{a}^{(V2)}$ (Superposición) & & \\ \hline
$V_{th} = V_{oc}$ & & \\ \hline
$R_{th}$ & & \\ \hline
$I_{N}$ & & \\ \hline
\end{tabularx}

\section*{Tabla 3: Registro de Resultados Integrados (Fases 3 a 6)}
Error Relativo: $\varepsilon_r = \left| \frac{\text{Medido} - \text{Teórico}}{\text{Teórico}} \right| \times 100\%$.

\vspace{0.3cm}
\noindent
\begin{tabularx}{\textwidth}{|>{\bfseries\RaggedRight}p{2.8cm}|>{\centering\arraybackslash}p{1cm}|>{\RaggedRight\arraybackslash}X|>{\RaggedRight\arraybackslash}X|>{\RaggedRight\arraybackslash}X|>{\RaggedRight\arraybackslash}X|>{\RaggedRight\arraybackslash}p{2.5cm}|}
\hline
\rowcolor{ucbblue!10}
\textbf{Caso} & \textbf{Var.} & \textbf{Teoría} & \textbf{SPICE} & \textbf{Medición} & \textbf{Error \%} & \textbf{Observación} \\ \hline
Superposición V1 & $V_{a}^{(V1)}$ & & & & & \\ \hline
Superposición V2 & $V_{a}^{(V2)}$ & & & & & \\ \hline
Red completa & $V_{a}$ & & & & & \\ \hline
Puerto abierto & $V_{oc}$ & & & & & \\ \hline
Equivalente & $R_{th}$ & & & & & \\ \hline
Norton & $I_{N}$ & & & N/A & N/A & Supervisada \\ \hline
Carga 1 ($1.0\,\text{k}\Omega$) & $V_L$ & & & & & \\ \hline
Carga 2 ($2.2\,\text{k}\Omega$) & $V_L$ & & & & & \\ \hline
Carga 3 ($4.7\,\text{k}\Omega$) & $V_L$ & & & & & \\ \hline
Carga 4 ($10\,\text{k}\Omega$) & $V_L$ & & & & & \\ \hline
\end{tabularx}

\end{document}
